\section{Comparing jump distributions and validating the estimator}
\label{sec:comparison}

\subsection{Ranking candidate jump laws}

With five interchangeable jump families, the applied question becomes:
which one does \emph{this} dataset support?
\code{JumpDistributionComparison} fits each candidate to the same
increments and reports two complementary kinds of evidence.

\paragraph{Information criteria.} AIC and BIC from \eqref{eq:aic-bic}
penalize the maximized likelihood for parameter count --- relevant here
because the candidates differ in dimension (Merton normal has one jump
parameter, the SGED four). They answer a \emph{relative} question: which
candidate trades fit against complexity best.

\paragraph{Goodness of fit.} Information criteria cannot say whether even
the best candidate actually fits. For that, the comparison computes a
Kolmogorov--Smirnov distance between the observed increments and the
fitted model's increment distribution --- the latter having no closed-form
CDF, it is represented by a large simulated reference sample from the
fitted model, so the statistic is a two-sample KS distance whose
reference-side noise is made negligible by sample size.

\begin{decision}[A bootstrap p-value, because parameters were estimated]
\label{dec:ks-bootstrap}
The classical KS p-value assumes the reference distribution was fixed
\emph{a priori}. Here it was estimated from the very data being tested,
which makes the naive p-value strongly anti-conservative --- the fitted
model always looks too good, the same phenomenon that motivates the
Lilliefors correction for Gaussian fits. The library restores validity
with the same device as Section~\ref{sec:jumptest}: simulate
$B$ datasets from the fitted model, \emph{re-fit the model on each}
(re-fitting is what accounts for estimation error), recompute each
replicate's KS distance to its own re-fitted model, and report the Monte
Carlo p-value \eqref{eq:mc-pvalue}. A model is inadequate when the
observed distance is large relative to the distances the model produces
against itself.
\end{decision}

\begin{lstlisting}[caption={Bootstrap KS p-value in \code{validation/distribution\_comparison.py} (excerpt).}]
for b in range(n_bootstrap):
    boot_data = self._simulate_increments(fitted_model, n_obs, data_seed)
    estimator = JumpDiffusionEstimator(
        boot_data, self.dt, jump_distribution=jump_distribution)
    rep_result = estimator.estimate(**estimate_kwargs)
    if not rep_result["convergence"]:
        continue
    model_b = JumpDiffusionModel(jump_distribution=jump_distribution,
                                 **rep_result["parameters"])
    reference_b = self._simulate_increments(model_b, reference_size, ref_seed)
    ks_b = float(ks_2samp(boot_data, reference_b).statistic)
    if ks_b >= ks_observed:
        exceed += 1
    n_successful += 1

return (1.0 + exceed) / (n_successful + 1.0)
\end{lstlisting}

\code{compare()} assembles everything into a table ranked by AIC (with
BIC, KS statistic and bootstrap p-value alongside), and
\code{plot\_comparison()} overlays each fitted mixture density on the
increment histogram --- the visual complement to the numbers. Because the
KS machinery relies only on simulation through
\code{JumpDistribution.rvs}, it works unchanged for any user-added jump
family, honoring the extension contract of Section~\ref{sec:contract}.

\subsection{Monte Carlo validation of the estimator}

Finally, \code{ValidationExperiment} closes the loop on the estimator
itself: choose true parameters, simulate $R$ independent paths, estimate
on each, and summarize the sampling behavior of the estimator ---
bias, standard deviation, RMSE per parameter, along with convergence
rates (all summary statistics condition on the runs whose optimization
converged). Two reporting conventions are deliberate. Relative errors are
reported as \code{NaN} whenever the true parameter is (numerically) zero
--- $\mu = 0$, \code{jump\_skew} $= 0$ and \code{jump\_loc} $= 0$ are all
natural experimental settings, and division by zero would otherwise
inject $\pm\infty$ into every downstream summary; \code{NaN} says
``undefined'', which is the honest answer. And the share of runs landing
within 5\% of the truth is reported as an \emph{accuracy rate}
(\code{within\_5pct\_rate}), pointedly not called ``coverage'': coverage
belongs to the confidence intervals of Section~\ref{sec:inference}, and
conflating the two invites exactly the kind of misreading a validation
tool exists to prevent.

This simulate--estimate--compare pattern is how every method in this
paper was tested during development, packaged for users; it is also the
seed of the full factorial Monte Carlo study (coverage of the three
interval routes, size and power of the jump test, identifiability maps)
that constitutes the research program around the library.
